% SPDX-License-Identifier: Apache-2.0
% covers: Ddr3Per
\datasheet{Ddr3Per}{The board's DDR3 memory as an AXI peripheral: AMD's MIG controller with its own AXI4 port, joined to the link by \code{AxiPerPins}, and making the design's clock.}

\dsfacts{\code{ddr3/src/lib.rs}}{\code{ddr3::Ddr3Per}}%
{\code{//docs:vreteno}, \code{//docs:pcie}}

\subsection*{Function}

\code{Ddr3Per} puts the DDR3 memory of the Alinx AX7A200B on an AXI
link. It is a unit of units with two children. \code{pins} is
\code{AxiPerPins<32, 32, 4, 5>} of \code{txhdl\_parts}, which puts the
link's peripheral end onto an AXI4 peripheral's pins and holds
nothing. \code{ctl} is \code{Ddr3}, AMD's MIG 7 Series controller,
which the build generates with its AXI4 port as
\code{//ddr3:ddr3\_mig\_axi}, behind the wrapper
\code{ddr3/hdl/ddr3\_axi32.v}. Thirty-seven AXI4 signals join the two.

There is no tracker in front of it. The controller keeps its own
transactions, cuts a burst into the memory's commands, and packs
thirty-two bit beats into its 256-bit native interface, so the
peripheral takes the link's channels as they are: \code{aw},
\code{ar} and \code{w} in, \code{b} and \code{r} out.

\code{Ddr3} is a foreign unit. Its netlist is an instance of the
module \code{ddr3\_axi32}, which the lowering names and does not write.
Its \code{run} is a model for simulation: \code{PinRam} of
\code{AxiPerPins} on the controller's pins, which holds the whole
gigabyte and is not ready until it has calibrated. The model does not
drive the memory's pins.

On the memory side its ports are the board's clock and the
controller's reset in, the design's clock and its reset out, since the
controller makes them, the \code{calib} output, and the memory's pins,
with \code{dq}, \code{dqs} and \code{dqs\_n} as pads both ways.

\subsection*{Parameters}

None. The controller calibrates in hardware, and its simulation
library provides the variant with the fast calibration, so neither
wants a parameter. The widths are the controller's: 32-bit data, five
identifier bits, and a 30-bit address, the memory's whole gigabyte.

\subsection*{Address map}

The link's address is 32 bits and the controller's 30, so the wrapper
drops the top two bits. On the board, \code{BoardRouter} sends the
peripheral every address that matches \code{0x4000\_0000} under the
mask \code{0xc000\_0000}, and dropping the two bits leaves the offset
from \code{0x4000\_0000}.

\dstables{Ddr3Per}

\subsection*{Behaviour}

\begin{itemize}
\item An address phase or a write beat at the head of its channel is
the controller's \code{valid} and fields, taken in the step its
\code{ready} is high. A response or a read beat the controller offers
goes onto its channel in the step the channel has room, and that room
is the controller's \code{ready}.
\item Every field of the link's address phase reaches a pin except the
region, which the controller has no pin for.
\item The wrapper gives the controller its AXI reset from the user
clock's reset, inverted and registered, never requests a refresh,
self-refresh or ZQ calibration, which the controller makes itself, and
gives it a constant temperature.
\item The controller makes the clocks. It takes the board's 200~MHz on
\code{sys\_clk}, which is also its delay reference, runs the memory
at 400~MHz, and hands out \code{ui\_clk}, the 100~MHz the design
runs on, with \code{ui\_rst} high until that clock is good. In the
netlist the controller's \code{i\_ui\_clk} is the design's clock,
which is that same clock fed back. \code{sys\_rst} is active high, a
pulse at power-on on the board.
\item \code{calib} is the controller's \code{o\_calib\_complete}.
\item In simulation the model is not ready for \code{MODEL\_WARMUP},
64 cycles. It offers a read's first beat \code{MODEL\_READ\_LATENCY},
24 steps, after it sees the address phase, and a write's response
\code{MODEL\_WRITE\_LATENCY}, 2 steps, after it sees the last beat. It
sees what the pins part drove a step after the pins part drove it.
Those are the controller's latencies as \code{//ddr3/sim:wrapper\_test}
measured them against the Micron models: a read's first beat 25
cycles after its address phase was taken, 27 for a burst of sixteen
and 63 right after a write, and a write's response 3 cycles after its
last beat.
\end{itemize}

\subsection*{Verification}

\begin{itemize}
\item \code{//ddr3:ddr3\_test} runs \code{words\_go\_in\_and\_come\_back}.
Behind an \code{AxiHost}, it writes words at \code{0x4000\_0000},
\code{0x4000\_0104} and \code{0x7fff\_fffc}, the first while the model
still calibrates. It reads them back, writes and reads back a burst of
sixteen, and checks that \code{calib} did not rise before the
warm-up.
\item \code{//ddr3:ddr3\_test} also runs
\code{the\_controller\_is\_instantiated\_and\_not\_written}. It checks
that the Verilog and the VHDL of \code{Ddr3Per} instantiate
\code{ddr3\_axi32}. The controller's clock must be on \code{clk}, the
board's clock and the clock it makes on the ports, and its pads on
the ports. The pins part's module must be written and the
controller's must not.
\item \code{//ddr3:ddr3\_test} runs the bandwidth tests of
\code{ddr3::bw}. With bursts in flight, reads take a word a cycle at
any latency up to 32 cycles, and one burst at a time pays the latency
once. Writes take a beat a cycle. The peripheral with its model
behaves as the pins part does at the model's latency.
\item \code{//cpu/vreteno:board\_test} runs
\code{the\_memory\_test\_runs\_on\_the\_board}, a compiled program that
writes and reads the memory through the model and must print
\code{ddr3 ok}. It also runs \code{the\_netlist\_holds\_the\_controller},
and \code{the\_ddr3\_path\_is\_timed\_by\_the\_core}, which finds a load
from the memory costing the model's read latency less two cycles more
than a load from the data memory.
\item \code{//ddr3/sim:wrapper\_test} simulates the wrapper and the
controller against the Micron models under Vivado's simulator. It
calibrates, then writes and reads back a burst of sixteen, a word in
another row, a word through an address with its top bits set, half a
word under its strobes, and a byte at an address that is not a
word's. It prints the latencies above. It is manual.
\code{//ddr3/sim:example\_test} is AMD's own proof of the controller as
configured, the native example design's traffic generator writing the
memory and reading it back through it; manual too.
\item \code{//cpu/vreteno/board/sim:}\allowbreak\code{board\_test} simulates the whole
board, this controller included, against two Micron models. It is
manual.
\end{itemize}

\subsection*{Used by}

\code{Board}, as \code{ddr3}, on the router's port for
\code{0x4000\_0000}, with no tracker in front.

\subsection*{Limits}

The Rust run of \code{Ddr3} is a model and not the controller. A
lowered netlist needs \code{//ddr3:hdl}, the wrapper, beside it, and
the controller's library, \code{//ddr3:ddr3\_mig\_axi}, which the build
generates under Vivado and which stays out of the tree, being AMD's.
The controller was generated without narrow bursts, so every beat is
the full thirty-two bits, with strobes for less; a byte's access is a
full-width beat at the byte's address under that byte's strobe, which
\code{wrapper\_test} checks. The path through the controller's own port
has run in simulation; its run on the board is the fourth step of
issue~\#1023.
